% Draft, 2026-10-28. All numbers come from results/ via paper/make_tables.py (tables in paper/tables/,
% numbers quoted in the text in paper/tables/stats.json). Re-run that script after any re-evaluation.
\documentclass[11pt]{article}

\usepackage[preprint]{acl}
\usepackage{times}
\usepackage{latexsym}
\usepackage[T1]{fontenc}
\usepackage[utf8]{inputenc}
\usepackage{microtype}
\usepackage{inconsolata}
\usepackage{booktabs}
\usepackage{multirow}
\usepackage{graphicx}

\newcommand{\todo}[1]{\textcolor{red}{[TODO: #1]}}

\title{Retrieved but Ignored: Where LLMs Lose Code Reuse When Refactoring}

\author{Evgenii Lialin \and Roman Vikhman \\
  \todo{affiliation} \\
  \texttt{\todo{emails}}}

\begin{document}
\maketitle

\begin{abstract}
When language models refactor code, they often re-implement logic that the repository already provides.
We ask where such reuse opportunities are lost: in retrieving the existing helper, in the model's decision
to call it, or in the edit itself. We build a benchmark of 129 test-validated reuse tasks from 12 Python
repositories by inlining a helper into one of its callers, so that the original caller is a known reference.
We then vary the context (none, BM25 snippets, the helper guaranteed among distractors, the whole module)
and the instruction (generic vs.\ explicit reuse) for four open models.
For three of the four models, the decision is the main bottleneck: even with the helper guaranteed in context
and a generic instruction, they ignore it in 62--81\% of tasks, while BM25 misses it in only 24\%.
An explicit reuse instruction raises reuse in every condition, by up to 2.8$\times$.
A newer coder model (Qwen3-Coder-30B-A3B) produces far more correct code than 7B models but does not reuse more;
only the newest general model (Qwen3.8-27B) does. Reuse also depends on where the helper appears
in the context and on its name: showing the helper under a neutral name lowers reuse substantially,
while direct memorisation of the original code is rare.
\end{abstract}

\section{Introduction}

Code assistants increasingly edit existing repositories rather than write code from scratch.
A well-known weakness is that they duplicate functionality the project already has:
AI-generated pull requests contain more redundant code than human ones \citep{huang2026codereuseinvestigatingcode},
and pretrained models frequently re-implement repository dependencies instead of calling them
\citep{hai-etal-2025-impacts}. Duplicated logic is a maintenance cost that tests do not reveal.

Refactoring benchmarks evaluate whether edits preserve behaviour or match human refactorings
\citep{xu2026swerefactorrepositorylevelbenchmarkrealworld,gautam2025refactorbenchevaluatingstatefulreasoning,thillen2026codetastellmsgeneratehumanlevel},
and retrieval-augmented systems supply repository context for completion, repair and refactoring
\citep{zhang-etal-2025-coderag,xiong-etal-2025-think,xu2025mantraenhancingautomatedmethodlevel}.
None of them isolates \emph{reuse}: whether a model, given a function that duplicates an existing helper,
replaces the duplicate with a call. When reuse fails, it is also unclear why: the helper may not have
been retrieved, the model may have seen it and ignored it, or the call may have been made incorrectly.

We study this with a controlled benchmark and a \emph{reuse funnel} that assigns every failure to exactly one
of these stages. Our contributions:
\begin{itemize}
  \setlength\itemsep{0pt}
  \item A benchmark of 129 reuse-refactoring tasks in 12 Python repositories, built by inlining a helper into its
        caller and kept only when the call site is covered by tests that still pass (\S\ref{sec:bench}).
  \item An evaluation of four open models under four context conditions and two instructions, with failures
        attributed to retrieval, decision or edit (\S\ref{sec:results}).
  \item Two controls that change the conclusions: the helper's position among retrieved snippets
        changes reuse for three of four models, and a renaming check separates naming cues from
        memorisation of the public repositories.
\end{itemize}

\section{Related Work}

\paragraph{Refactoring with LLMs.}
SWE-Refactor \citep{xu2026swerefactorrepositorylevelbenchmarkrealworld} and MANTRA
\citep{xu2025mantraenhancingautomatedmethodlevel} evaluate behaviour-preserving refactorings in Java;
RefactorBench \citep{gautam2025refactorbenchevaluatingstatefulreasoning} poses multi-file Python refactorings
to agents; CodeTaste \citep{thillen2026codetastellmsgeneratehumanlevel} asks whether models choose
human-like refactorings. They measure correctness or agreement with humans, not reuse of existing code.

\paragraph{Repository context and reuse.}
RepoExec \citep{hai-etal-2025-impacts} introduces the dependency invocation rate for repository-level
generation and finds that models often re-implement available dependencies; DependEval
\citep{du-etal-2025-dependeval} benchmarks dependency understanding. We adopt the idea of measuring
invocation, but for refactoring, where the duplicated logic is already present and must be recognised.

\paragraph{Retrieval for code.}
CodeRAG \citep{zhang-etal-2025-coderag}, ReflectCode \citep{liu-etal-2025-beyond-function},
CoRet \citep{fehr-etal-2025-coret} and Think-Search-Patch \citep{xiong-etal-2025-think} improve retrieval
for completion, change requests, editing and repair. We use BM25 \citep{robertson2009bm25} and an oracle
to separate retrieval quality from what the model does with the retrieved code.

\paragraph{Instructions and context position.}
CodeIF \citep{yan-etal-2025-codeif} shows that instruction following in code generation varies widely;
we compare a generic refactoring instruction with an explicit reuse instruction. Models use information
unevenly across long contexts \citep{liu-etal-2024-lost}; we test whether the helper's position among
retrieved snippets matters.

\section{Benchmark}
\label{sec:bench}

\paragraph{Task construction.}
For each repository we find statement-level calls from a function $F$ to a repository helper $H$
(a module function, a method called via \texttt{self}, or a function of an imported module).
We inline $H$'s body at the call site, binding arguments, renaming locals to avoid clashes and
turning early returns of $H$ into returns of $F$. The mutated $F$ duplicates $H$; the original $F$,
which calls $H$, is the reference. A task is kept only if the call site is covered by the test suite
and the full suite still passes after inlining. Inlining an overridable method would change subclass
behaviour, so methods defined by more than one class are skipped.

\paragraph{Data.}
We use 12 Python libraries with fast test suites (boltons, cachetools, click, mistune, parso, pycparser,
pyflakes, pygments, python-tabulate, sqlparse, tomlkit, toolz) at pinned commits. We keep every task whose
helper lives in another module (36 tasks) and a seeded sample of same-module tasks, up to 12 per repository,
giving 129 tasks. Helpers have 2--30 lines (median 7); 64 are methods called via \texttt{self}, 61 plain
functions and 4 module attributes.

\paragraph{Conditions.}
All conditions except \textsc{file} share a base context: the module's imports and the mutated $F$
(with its class header). \textsc{none} adds nothing. \textsc{bm25} adds the top-5 repository functions
retrieved with BM25 using $F$ as the query; the original $F$ is excluded from the index.
\textsc{oracle} adds $H$ and the top-4 BM25 distractors, with $H$ at a seeded random position;
\textsc{oracle-first}/\textsc{last} place it first or last (\S\ref{sec:position}).
\textsc{file} shows the whole mutated module, truncated to 30k characters around $F$ (53 tasks);
it contains $H$ only for same-module helpers (80 tasks after truncation).
Each condition is crossed with two instructions: \emph{generic} (``Refactor this function to improve its
readability and maintainability'') and \emph{reuse}, which adds ``Where possible, reuse existing functions
from the project instead of re-implementing their logic.'' Both ask the model to keep the name, signature
and behaviour and to output only the changed function and any new helpers.

\paragraph{Evaluation.}
The model's function replaces $F$; new helper functions or methods are added, but redefinitions of existing
code are dropped and the output is marked non-compliant. We run the repository's tests in a sandbox.
An output is \emph{valid} if it is compliant and the tests pass; it \emph{reuses} the helper if the new $F$
calls $H$ or a near-identical copy of $H$ elsewhere in the repository, resolved through imports.
The \emph{reuse funnel} assigns each task to one stage: $H$ \emph{not in context}, in context but
\emph{ignored}, called but \emph{broken} (tests fail), or success. Conditions and models are compared
per task with exact McNemar tests.

\section{Experimental Setup}

We evaluate Qwen2.5-Coder-7B-Instruct \citep[Q2.5C-7B;][]{hui2024qwen25codertechnicalreport} and
DeepSeek-Coder-6.7B-Instruct \citep[DSC-6.7B;][]{guo2024deepseekcoderlargelanguagemodel} in bf16 with
Hugging Face Transformers, and two 2025--26 models in 4-bit quantisation with vLLM
\citep{kwon2023efficientmemorymanagementlarge}: Qwen3-Coder-30B-A3B-Instruct
\citep[Q3C-30B, a mixture of experts with 3B active parameters;][]{yang2025qwen3technicalreport,qwen3coder-awq}
and Qwen3.8-27B \citep[Q3.8-27B, thinking disabled;][]{qwen38-int4}. All runs use greedy decoding,
at most 2{,}048 new tokens, and a single 24\,GB GPU; each model answers 1{,}032 prompts
(129 tasks $\times$ 4 conditions $\times$ 2 instructions), plus 516 for the position ablation.

\section{Results}
\label{sec:results}

\begin{table}[t]
\centering
\small
\setlength{\tabcolsep}{3.5pt}
\begin{tabular}{llcccc}
\toprule
Model & Instr. & \textsc{none} & \textsc{bm25} & \textsc{oracle} & \textsc{file} \\
\midrule
Q2.5C-7B & generic & 1 \small(49) & 21 \small(60) & 38 \small(59) & 18 \small(52) \\
 & reuse & 3 \small(43) & 43 \small(56) & 53 \small(53) & 22 \small(52) \\
\midrule
DSC-6.7B & generic & 4 \small(36) & 33 \small(45) & 36 \small(42) & 18 \small(47) \\
 & reuse & 5 \small(37) & 47 \small(45) & 45 \small(42) & 19 \small(43) \\
\midrule
Q3C-30B & generic & 2 \small(91) & 22 \small(91) & 19 \small(88) & 9 \small(93) \\
 & reuse & 5 \small(90) & 41 \small(87) & 54 \small(83) & 27 \small(88) \\
\midrule
Q3.8-27B & generic & 5 \small(90) & 50 \small(92) & 53 \small(88) & 36 \small(96) \\
 & reuse & 12 \small(92) & 74 \small(92) & 84 \small(88) & 64 \small(93) \\


\end{tabular}
\caption{Helper reuse in \% of 129 tasks (valid outputs in \% in parentheses).}
\label{tab:main}
\end{table}

\begin{table}[t]
\centering
\small
\setlength{\tabcolsep}{3pt}
\resizebox{\columnwidth}{!}{%
\begin{tabular}{lcccc}
\toprule
& \multicolumn{2}{c}{\textsc{bm25}} & \multicolumn{2}{c}{\textsc{oracle}} \\
Model & generic & reuse & generic & reuse \\
\midrule
Q2.5C-7B & 31/73/4/21 & 31/49/9/40 & 0/80/14/35 & 0/60/17/52 \\
DSC-6.7B & 31/60/13/25 & 31/41/21/36 & 0/83/19/27 & 0/71/21/37 \\
Q3C-30B & 31/73/0/25 & 31/50/4/44 & 0/105/0/24 & 0/59/4/66 \\
Q3.8-27B & 31/40/1/57 & 31/13/3/82 & 0/60/5/64 & 0/21/3/105 \\

\end{tabular}}
\caption{Reuse funnel over 129 tasks: not in context / ignored / broken / success.}
\label{tab:funnel}
\end{table}

\paragraph{Without context, models re-implement.}
In \textsc{none}, reuse is 1--5\% for three models and at most 12\% for Q3.8-27B (Table~\ref{tab:main}).
Instead, the smaller models write a new helper of their own in 68\% (Q2.5C-7B) and 48\% (DSC-6.7B)
of outputs (Q3C-30B 15\%, Q3.8-27B 10\%).
Many of the remaining calls without context are explained by the prompt itself: the helper's name already
appears in the module's imports in 44 of the 258 \textsc{none} prompts.

\paragraph{The decision, not retrieval, is the main loss.}
BM25 places the helper in the top 5 for 76\% of tasks (MRR 0.58), so retrieval loses 31 tasks.
With the helper guaranteed in context and the generic instruction, three models ignore it in 80--105
of 129 tasks, and even Q3.8-27B in 60 (Table~\ref{tab:funnel}). Calls that break the tests are rare
for the newer models (0--5) but not for the 7B models (up to 21).
Only for Q3.8-27B with the reuse instruction does retrieval become the larger loss (31 not retrieved vs.\ 13 ignored).
Whole-module context yields the least reuse of the three context conditions for every model,
although the helper is visible in 80 tasks.

\paragraph{An explicit instruction raises reuse.}
Adding the reuse sentence raises reuse in every condition with context, e.g.\ \textsc{oracle} from 19\% to 54\%
for Q3C-30B and from 53\% to 84\% for Q3.8-27B, while validity drops by at most 6 points.
Models can reuse, but they do not consider it part of refactoring unless asked.

\paragraph{Stronger is not the same as more reusing.}
Q3C-30B produces valid code in 83--93\% of cases versus 36--60\% for the 7B models, yet it does not
reuse significantly more than Q2.5C-7B in any condition; with \textsc{oracle} and the generic instruction it reuses less
(10 vs.\ 35 discordant tasks, $p<.001$). Q3.8-27B reuses more than both in every condition with context
($p\le.002$), e.g.\ \textsc{bm25} with the reuse instruction 74\% vs.\ 41--47\%.

\subsection{Position of the helper}
\label{sec:position}

\begin{table*}[t]
\centering
\small
\begin{tabular}{lccccc|ccccc}
\toprule
& \multicolumn{5}{c|}{generic} & \multicolumn{5}{c}{reuse} \\
Model & first & random & last & first/last only & $p$ & first & random & last & first/last only & $p$ \\
\midrule
Q2.5C-7B & 40 & 49 & 41 & 13/14 & 1.0 & 63 & 69 & 67 & 12/16 & .572 \\
DSC-6.7B & 47 & 46 & 47 & 11/11 & 1.0 & 75 & 58 & 55 & 25/5 & <.001 \\
Q3C-30B & 37 & 24 & 25 & 14/2 & .004 & 70 & 70 & 66 & 17/13 & .585 \\
Q3.8-27B & 78 & 69 & 69 & 12/3 & .035 & 110 & 108 & 108 & 4/2 & .688 \\

\end{tabular}
\caption{Tasks (of 129) where the helper is reused when placed first, at a seeded random position, or last among
the five \textsc{oracle} snippets; discordant tasks and exact McNemar $p$ for first vs.\ last.}
\label{tab:position}
\end{table*}

DSC-6.7B reused the helper more with \textsc{bm25} than with \textsc{oracle}, although \textsc{oracle} always
contains it. BM25 usually ranks the helper first, so we repeated \textsc{oracle} with the helper first or last,
with identical snippets otherwise (Table~\ref{tab:position}). Three models show a primacy effect, but under
different instructions: DSC-6.7B only with the reuse instruction (75 vs.\ 55 tasks, $p<.001$), Q3C-30B and
Q3.8-27B only with the generic one (37 vs.\ 25 and 78 vs.\ 69). Q2.5C-7B shows none.
Reordering identical snippets changes the outcome of up to 30 tasks (23\%) per model and instruction. Oracle upper
bounds should therefore fix the helper's position, and differences between conditions should be tested per
task rather than read from aggregate rates.

\subsection{Memorisation or naming?}
\label{sec:renaming}

\begin{table}[t]
\centering
\small
\setlength{\tabcolsep}{3pt}
\resizebox{\columnwidth}{!}{%
\begin{tabular}{llccc}
\toprule
Model & Instr. & \textsc{bm25} & \textsc{oracle} & \textsc{file} \\
\midrule
Q3C-30B & generic & 23$\to$9\,\small(<.001) & 19$\to$12\,\small(.065) & 11$\to$8\,\small(.508) \\
 & reuse & 46$\to$36\,\small(.031) & 59$\to$46\,\small(.004) & 29$\to$22\,\small(.143) \\
\midrule
Q3.8-27B & generic & 56$\to$22\,\small(<.001) & 61$\to$33\,\small(<.001) & 39$\to$21\,\small(<.001) \\
 & reuse & 83$\to$62\,\small(<.001) & 96$\to$82\,\small(<.001) & 70$\to$59\,\small(.003) \\
\bottomrule

\end{tabular}}
\caption{Helper reused (of 114 tasks) with the original name $\to$ with a neutral name in the prompt
(McNemar $p$).}
\label{tab:renamed}
\end{table}

All repositories are public and predate the models, so a model might reuse the helper because it remembers
the original caller. We re-ran the two newer models with the helper shown under a neutral name
(\texttt{fn\_<hash>}) everywhere in the prompt; outputs are mapped back before testing, and any use of the
original, never-shown name is counted as memorised and broken. We use the 114 tasks whose helper name is
unique and not also a builtin or standard-library name.
Memorisation is rare: 18 of 912 outputs for Q3.8-27B and 5 for Q3C-30B use the original name, almost all
for private sibling methods such as \texttt{self.\_parse\_key}. Without context, reuse falls to 2 of 114
for both models. Neutral names, however, reduce reuse substantially for both models (Table~\ref{tab:renamed}):
with \textsc{bm25} and the generic instruction both lose about 60\% of their reuse. Since the control model
loses a similar share, the drop reflects the loss of a descriptive name as a cue rather than memorisation
specific to Q3.8-27B. Q3.8-27B still reuses far more than Q3C-30B under renaming
(\textsc{oracle}, reuse instruction: 37 vs.\ 1 discordant tasks, $p<.001$).

\section{Conclusion}

Reuse in LLM refactoring is lost mainly at the model's decision, not at retrieval: models see the helper and
ignore it, unless explicitly asked to reuse. Better retrieval alone would therefore recover few tasks for most
models. Higher code quality does not imply more reuse, and reuse depends on surface cues such as the helper's
position and name. Benchmarks of repository-aware code editing should report reuse alongside test results,
fix the position of oracle context and control for naming. We release the benchmark, prompts, outputs and
evaluation code.\todo{repository link}

\section*{Limitations}

Tasks are created by deterministic inlining, so the duplicated code is a near-verbatim copy of the helper;
real duplication is usually paraphrased, which would make reuse harder to recognise, and a set of tasks
mined from real commits is future work. We use 12 popular Python libraries; results may differ for other
languages and for private code. Passing tests do not prove equivalence, although we only keep tasks whose
call site is covered. We evaluate four open models with greedy decoding and a single run, two of them
quantised to 4 bits; we do not evaluate proprietary models, dense or agentic retrieval. Our renaming check
changes names only in the prompt; it does not rule out memorisation of code structure. A small number of
outputs were re-evaluated on a different machine after harness fixes, where test outcomes can differ slightly.

\bibliography{references}

\end{document}
